\section{Further questions}
\label{sec:questions}

The first question asks whether binary moves connect every represented
fiber on the full odd-prime support.

\begin{question}\label{question:full-support-connectivity}
For the full binary Goldbach ring $G$ on the odd primes, is every
nonempty fiber $\mathcal M_{n,d}$ connected by binary moves?
Equivalently, is the natural map
\[
 G\longrightarrow\kk[u^pv:p\text{ an odd prime}]
\]
an isomorphism, and is $G$ a domain?
\end{question}
The equivalence with the toric presentation follows from the component
basis in Proposition~\ref{setup:components}. For the full support,
Theorem~\ref{depth:prime-pair-group} identifies any two monomials of
the same bidegree after multiplication by a common monomial.
Such a multiplier is nonzero in $G$. Thus failure of connectivity
would give a nonzero zero divisor, while the toric quotient is a
domain. The localization argument in Appendix~\ref{app:saturation}
does not by itself permit cancellation before localization.

For each fixed monomial degree $d\ge3$, Theorem~\ref{main:weightwise} gives a
positive lower bound for the fraction of the linear strand spanned by
cycles from positive common shifts. Already in monomial degree three there is
a finer question.

\begin{question}
For the full binary ring on the odd primes, let $W^+_{n,3}$ be the span
of the positive common-shift classes at weight $n$. For odd $n$ tending
to infinity, does
\[
 \frac{\dim_\kk W^+_{n,3}}{\beta_{2,n,3}}\longrightarrow1?
\]
How does the answer depend on the coefficient field?
\end{question}

\begin{question}
For fixed $i$, what is the growth of
$\beta^{<m}_{i,*,*}$ and $\beta^{\le m}_{i,*,*}$
when all monomial degrees are included?
\end{question}
Theorem~\ref{grow:both-cutoffs} determines the linear contribution.
For the binary variable cutoff the next strand has lower order.
Estimating the total requires control of syzygies in all later
monomial degrees. One may also let $i$ vary with $m$, a regime
outside the asymptotic formula for fixed homological degree.

\begin{question}
Does $\depth G^{<m}\le2$ hold at every variable cutoff? Which
cutoffs have depth zero, one, or two? Is the depth of the binary
relation-weight truncations uniformly bounded?
\end{question}
Theorem~\ref{depth:prime-pair-group} identifies the stable integral
group with $\mathbb Z^2$. Corollary~\ref{depth:bounded-depth} gives
the bound two for all sufficiently large variable cutoffs, while
Corollary~\ref{depth:verified-range} verifies it for
$3\le m\le20\,000\,000$. We do not know that these two ranges
overlap. The exact depth at each cutoff and the universal bound
remain open. The prime-pair condition at $3$ is sufficient for the
bound, but the depth question does not assert that this condition
is necessary.
For $G^{\le m}$ the bound in Theorem~\ref{depth:relation-cutoff-pd}
uses only pair relations whose weight is at most $m$.
Conjecture~\ref{depth:dual-conjecture} proposes a nonzero infinite
dual Koszul class. Its finite representatives give explicit depth
bounds at the checked cutoffs; Proposition~\ref{depth:dual-prime-pair}
relates the conjecture to the prime-pair condition for $x_3$-saturation.

\subsection*{Acknowledgments}
The author thanks his Clemson University Ph.D. advisor,
Keri Sather-Wagstaff, for frequent conversations and meetings and for
encouraging him to pursue this research.

\subsection*{Computational data}
The arithmetic supplement accompanying this paper contains the
illustrated resolution strands, the weight-$33$ Koszul cycles, the
linear-strand constants, and the saturation, variable-quotient, and
dual-cocycle calculations in Appendix~\ref{app:saturation}. It also
contains the exhaustive prime-pair check for
Corollary~\ref{depth:verified-range}. The supplement is supplied in
\path{anc/supplement} with the arXiv sources; paths in the text are
relative to its parent directory.
Its \path{README.txt} identifies the assertions checked, dependencies,
and expected output. The \path{verify.py} runner verifies the manifest
hashes and executes six standard-library Python checks, including the
prime-pair sieve through $20\,000\,000$. The ideal, exact-depth, and
Tor calculations in Appendix~\ref{app:saturation} require running the
three supplied \textsc{Singular} scripts separately; the Python runner
does not execute those scripts. The finite depth corollary uses the
prime-pair check and the localization argument in its proof, not a
computed minimal resolution. The arithmetic asymptotic and
interpolation theorems are proved in the text.

\Needspace{17\baselineskip}
\subsection*{Use of generative AI}
The Goldbach-ring framework, the construction and existence proof of its
minimal bigraded free resolution, the study of cellular models for its
fixed-bidegree complexes, and the shifted-Koszul-cycle proof of
infinite projective dimension originated with the author. The initial
computational identification of characteristic-two dependence, treated
in the companion paper, also originated with the author.

During the subsequent development of this work, the author used OpenAI's
GPT-6 Astra for substantive mathematical assistance. The model proposed
the results stated as Theorems~\ref{main:weightwise},
\ref{main:cutoff-growth}, and~\ref{main:prime-interpolation} and supplied
crucial additive-number-theoretic arguments and initial proof drafts,
building on the author's constructions. The author directed this
development, checked and verified the arguments, and revised their
presentation. AI assistance was also used to develop computational code
and verification routines and to organize and edit the manuscript.

The author takes responsibility for the mathematical claims, proofs,
computations, references, and final presentation.
